% Statements shared by the TISMIR and arXiv versions.
\newcommand{\ackText}{This work was carried out in connection with the Spr\aa kbanken project at the Division of
Speech, Music and Hearing, KTH Royal Institute of Technology.}

\newcommand{\contribText}{HC conceived the project and its goal of an endless guqin ``AI radio'', and directed every
stage of the exploration. HC screened the library by listening to whole recordings, decided which performers and
pieces to keep or exclude (including recordings with noise floors, ensembles, applause or voice), and noticed that a
DVD collection might duplicate the CD material. HC required that splits be made by composition rather than by
performer, that evaluation openings never repeat a piece or performer, and that validation recordings stay out of
training so that overfitting could be measured. HC proposed moving from text-free continuation to mood labels
derived from piece titles, and then to captions telling the story behind each title; proposed the move towards a
language-model-like generator over latents; observed that idiomatic guqin semitones occur in slides rather than at
the pluck, which shaped the onset-based measure; proposed random crop lengths in addition to random start points;
asked for training to continue until validation loss genuinely stopped improving, and for the run to be resumable
after power cuts. HC performed all listening studies and wrote all ratings and comments, identified the silent
seams and the loudness drift in chained continuation, and selected the takes used for demonstration.}

\newcommand{\aiText}{AI coding assistants (OpenAI Codex and Anthropic Claude) implemented and ran the code for data
processing, training, generation, listening pages and analysis under the author's direction, searched for and
summarised background sources on the pieces, and drafted parts of this manuscript. All research decisions and all
listening judgements are the author's, who reviewed the results and the text and takes responsibility for the
content.}
